{{دیوان
| عنوان = ایک شام
| ماخذ = ویکی ماخذ
| ماخذ_لنک = https://ur.wikisource.org/wiki/%D8%A7%DB%8C%DA%A9_%D8%B4%D8%A7%D9%85
| مقابلہ = اقبال سائبر لائبریری (اقبال اکادمی پاکستان)
| مقابلہ_لنک = https://www.iqbalcyberlibrary.net
}}
=== (دریائے نیکر، ہائیڈل برگ کے کنارے پر) ===

<poem>
خاموش ہے چاندنی قمر کی
شاخیں ہیں خموش ہر شجر کی

وادی کے نوا فروش خاموش
کُہسار کے سبز پوش خاموش

فطرت بے ہوش ہو گئی ہے
آغوش میں شب کے سو گئی ہے

کچھ ایسا سکُوت کا فُسوں ہے
نیکر کا خرام بھی سکُوں ہے

تاروں کا خموش کارواں ہے
یہ قافلہ بے درا رواں ہے

خاموش ہیں کوہ و دشت و دریا
قُدرت ہے مُراقبے میں گویا

اے دِل! تُو بھی خموش ہو جا
آغوش میں غم کو لے کے سو جا
</poem>
